\section{Open questions}
\label{sec:conclusion}

The main approximation question for treedepth remains the gap
between the $O(\log n)$ upper bound~\cite{Anand2026} and
constant-factor hardness. Our SSE result excludes every fixed
constant factor, while leaving an $o(\log n)$ approximation
possible. A superconstant hardness bound would further clarify
this gap.

Several hardness thresholds remain open. For treedepth, can the
NP-hardness factor exceed
$18/17$, or can ordinary UGC exclude all fixed constant factors?
For treewidth, the corresponding thresholds proved here are
$12/11$ and $4/3$, whereas all fixed factors are already excluded
under SSEH~\cite{APW}.

For total elimination depth, the $8/7$ ordinary-UGC lower bound
remains far from the $O(\sqrt{\log n})$ approximation inherited
from average-case graph search~\cite{Szyfelbein2025GraphSearch}.
The NP-hardness construction gives an absolute factor above one.
Obtaining a larger NP-hardness gap is a separate question from
improving the UGC bound.

Finally, our degree bounds depend on the excluded approximation
ratio. Preserving comparable gaps on graphs of a prescribed small
maximum degree would sharpen the distinction between hardness of
computing the optimum and approximation hardness in sparse graphs.
